\medskip

Un vendeur de voitures possède un stock de \np{1000} voitures dont les caractéristiques sont résumées dans le tableau ci-dessous.

\medskip

\begin{tabularx}{\linewidth}{|l|*{4}{>{\centering \arraybackslash}X|}}\hline
&Blanche &Noire &Rouge&TOTAL\\ \hline
Française &150 &$x$&400&750\\ \hline
Étrangère &100& 50 &100& 250\\ \hline
TOTAL &250&250 &500&\np{1000}\\ \hline
\end{tabularx}

\begin{enumerate}
\item 
Le nombre $x$ représente le nombre de voitures françaises noires.

Il y a 250 voitures noires, dont 50 étrangères; donc $x=250-50=200$.

\item 
Il y a \np{1000} voitures dans le stock dont 250 voitures noires.

Le pourcentage de voitures noires parmi les voitures du stock est donc égal à\\
 $\dfrac{250}{\np{1000}}$ c'est-à-dire $\dfrac{25}{100}$ soit 25\,\%.

\item 
Il y a \np{1000} voitures dans le stock dont 50 voitures noires étrangères.

Le pourcentage de voitures noires étrangères parmi les voitures du stock est donc égal à
 $\dfrac{50}{\np{1000}}$ c'est-à-dire $\dfrac{5}{100}$ soit 5\,\%.

\item 
Il y a 750 voitures françaises dont 150 voitures blanches.

Le pourcentage de voitures blanches parmi les voitures françaises est donc égal à\\
 $\dfrac{150}{\np{750}}$ c'est-à-dire $\dfrac{1}{5}$ ou $\dfrac{20}{100}$ soit 20\,\%.

\item 
Il y a 250 voitures blanches dont 150 voitures françaises.

Le pourcentage de voitures françaises parmi les voitures blanches  est donc égal à\\
 $\dfrac{150}{\np{250}}$ c'est-à-dire $\dfrac{15}{25}$ ou $\dfrac{60}{100}$ soit 60\,\%.

\item 
\begin{list}{\textbullet}{ Alice et Benoît jouent au jeu suivant:}
\item Alice choisit au hasard une voiture parmi les voitures Françaises. \\
Elle remporte 1 euro si ce n'est pas une voiture rouge.

Parmi les 750 voitures françaises, il y en a $750-400=350$ qui ne sont pas rouges. La probabilité pour Alice de gagner est donc $\dfrac{350}{750}$.

\item Benoit choisit au hasard une voiture parmi les voitures Blanches. \\
Il remporte 1 euro si c'est une voiture étrangère.

Parmi les 250 voitures blanches, il y en a 100 qui sont étrangères. La probabilité pour Benoit de gagner est donc $\dfrac{100}{250}$.
\end{list}

Or $\dfrac{100}{250} = \dfrac{300}{750} < \dfrac{350}{750}$,
donc c'est Alice qui a le plus de chance de remporter 1 euro.
\end{enumerate}

\bigskip
